\chapter{Implementation and Validation}

\section{Introduction}

JANUS is implemented in Python and PowerShell around an existing Wazuh
single-node laboratory. FastAPI exposes the control surface, Streamlit provides
the dashboard, SQLite stores the registry and evidence, and Docker runs the
Wazuh manager, Indexer, and dashboard. The current validated host is Windows.

\section{Project structure}

\begin{lstlisting}[caption={Simplified JANUS source tree},label={lst:tree}]
HoneyDocJANUS/
|-- main.py
|-- src/
|   |-- alerting/        # FastAPI server and Streamlit dashboard
|   |-- core/            # configuration, database, evidence, safety
|   |-- janus/           # DOCX and multi-format generators
|   |-- janus_v2/        # registry, Wazuh, telemetry, verdicts
|   |-- lifecycle/       # atomic deployment and TTL rotation
|   `-- tactical/        # prompt, LLM, coherence and scoring
|-- scripts/             # setup, Wazuh rules, tests, start/stop
|-- corpus/              # controlled financial, HR and technical context
|-- data/shared/         # monitored deployment root (ignored by Git)
|-- docs/
`-- tests/
\end{lstlisting}

\section{Core implementation}

\subsection{Startup and configuration}

\texttt{main.py} initializes directories and the database, starts the FastAPI
application, and conditionally starts the rotation manager and experimental
components. Fake SSH/HTTP services and experimental prompt/credential layers
remain disabled. On application startup, the Wazuh detection collector and the
forensic collector are created only when their configuration is complete.

Configuration is loaded from \texttt{.env}; the real file is ignored by Git.
The API key protects administrative routes. Validation functions reject invalid
intervals, missing provider configuration in fail-closed mode, and unsafe path
settings.

\subsection{AI content and safety gate}

The tactical layer builds a business prompt from a controlled corpus. If a
remote LLM key is configured, the model returns a narrative; otherwise a local
fallback keeps the pipeline operational. The output passes a coherence check
and forbidden-pattern scan. Token allocation happens only after this gate,
which avoids wasting tokens or persisting external identifiers for rejected
content.

\subsection{DOCX assembly}

The DOCX generator writes visible business content and inserts a non-macro
Office beacon in the document footer. The OOXML contains one
\texttt{INCLUDEPICTURE} field and one external image relationship whose target
is the Word token URL. The final inspector opens the DOCX as a ZIP package and
requires exactly one matching field and relationship, no VBA parts, and no
visible JANUS fingerprint.

\subsection{XLSX and structured formats}

The accounting workbook generator creates balanced journal, income statement,
balance-sheet, and budget data. The workbook uses supported table/filter
structures and an external token reference that opens without Excel repair.
CSV, JSON, YAML, ENV, and ZIP generators create format-valid content and embed
only the token material allowed by policy. For cloud credentials, the provider's
exact decoy key pair is used; fallback fake fields are not invented.

\subsection{Deployment and registry}

The path policy resolves a requested target and rejects traversal outside
\texttt{JANUS\_DEPLOY\_ROOT}. The injector writes the artifact atomically,
computes SHA-256, and persists the document and token transaction. Generated
files, the database, logs, backups, and \texttt{.env} are excluded from Git.

\section{Wazuh and evidence processing}

PowerShell setup scripts configure the Windows audit subcategories, inherited
SACL, Wazuh manager rules, and Indexer collector settings. The collector queries
the configured rule set, handles overlap and search-after cursors, normalizes
4663 access masks into read/modify/delete, matches paths against active decoys,
and applies a deployment-grace suppression policy.

Every accepted alert is stored with Wazuh identifiers, the normalized action,
process, user, agent, verdict, raw JSON, and a canonical raw-evidence SHA-256.
The Wazuh rule identifier is also normalized into a top-level field so the
dashboard can display rules such as 100101 directly.

\section{Dashboard}

The Alert page intentionally presents two independent sections: local Wazuh
detections and Canarytokens callbacks received by JANUS. In the public
email-only configuration, the interface explains that provider emails are not
copied into the local callback table. A non-blocking Streamlit fragment refreshes
the alert view every five seconds without interfering with navigation.

The HoneyDocs page displays safe metadata only. The Generate page keeps the
dependent type and format selectors outside the submission form, preventing a
stale DOCX value after choosing JSON. It also explains the activation condition
before and after generation.

\section{Automated validation}

The final suite contains 98 tests and 7 subtests. It covers:

\begin{itemize}
  \item API authentication and callback rejection;
  \item path traversal and deployment-root enforcement;
  \item DOCX/XLSX beacon structure and absence of macros;
  \item structured format validity and ZIP path safety;
  \item Canarytokens provider selection and secret handling;
  \item SQLite migrations, registry, raw-evidence hashes, and cursor behavior;
  \item Wazuh parsing, collection, suppression, verdicts, and timelines;
  \item dashboard navigation and non-blocking refresh;
  \item local LLM fallback and coherence behavior.
\end{itemize}

The validated output is:

\begin{lstlisting}[caption={Automated acceptance result},label={lst:pytest}]
98 passed, 1 warning, 7 subtests passed
\end{lstlisting}

The remaining warning concerns a Starlette/httpx test-client deprecation and
does not represent a functional failure.

A final endurance check exposed contention between the Wazuh collector and the
forensic-retention transaction in the shared SQLite database. The collector
stores now use a common re-entrant writer lock, a 60-second busy timeout, and an
immediate retention transaction. A concurrent-writer regression test was added.
After an offline backup, bounded pruning, and vacuum, the database passed
\texttt{quick\_check}, retained all 130 Wazuh detections, and both collectors
completed successive polls with zero errors.

\section{Live Windows and Wazuh validation}

On 10 August 2026, the stack was restarted after a host reboot. Docker reported
the Wazuh manager, Indexer, and dashboard as running. The JANUS API, dashboard,
Wazuh collector, forensic collector, and Canarytokens provider all reported a
healthy state.

A controlled binary read of HoneyDoc 26 generated Windows Security event 4663.
Wazuh applied rule 100101, described as a JANUS read/open candidate. The JANUS
collector correlated the event with the active document and stored:

\begin{itemize}
  \item the audited object path and HoneyDoc identifier;
  \item action \texttt{read} and access mask \texttt{0x1};
  \item the PowerShell process and local test account;
  \item Wazuh rule identifier 100101;
  \item raw evidence and its SHA-256 hash;
  \item a low verdict because no expected host or expected user was configured.
\end{itemize}

\begin{figure}[H]
\centering
\includegraphics[width=0.95\linewidth]{wazuh-rule-100101.png}
\caption{Wazuh document details showing JANUS rule 100101}
\label{fig:wazuh-rule}
\end{figure}

\begin{figure}[H]
\centering
\includegraphics[width=0.95\linewidth]{wazuh-4663-fields.png}
\caption{Event 4663 fields: object, process, access mask, and account}
\label{fig:wazuh-fields}
\end{figure}

\begin{figure}[H]
\centering
\includegraphics[width=0.95\linewidth]{windows-sacl-4663-test.png}
\caption{Inherited SACL and controlled Windows 4663 test}
\label{fig:sacl-test}
\end{figure}

\begin{figure}[H]
\centering
\includegraphics[width=0.95\linewidth]{janus-telemetry-status.png}
\caption{JANUS forensic telemetry status endpoint}
\label{fig:telemetry-status}
\end{figure}

\section{Multi-format and Canarytokens validation}

Manual tests produced valid DOCX, XLSX, JSON, and AWS ENV artifacts. The JSON
selector defect was corrected and verified through the dashboard. An offline
inspection of a newly generated DOCX confirmed one external Office relationship
and no macro. Excel opened the corrected workbook without repair and generated
a Canarytokens email containing a public source IP, provider geolocation, and
the user-agent \texttt{Mozilla/4.0 (compatible; ms-office; MSOffice 16)}. The
exact IP must be redacted in a public report.

The AWS ENV test behaved differently by design. Opening the file produced a
local Wazuh signal but no remote email. The official token is expected to alert
only when the decoy keys are presented to the AWS API.

\begin{table}[H]
\centering
\caption{Observed and expected acceptance results}
\label{tab:acceptance-results}
\small
\begin{tabular}{p{0.22\linewidth}p{0.27\linewidth}p{0.27\linewidth}p{0.16\linewidth}}
\toprule
Scenario & Local Wazuh layer & Remote token layer & Result \\
\midrule
DOCX read & 4663 and read/open rule & Conditional on Word policy & Local validated; remote conditional \\
XLSX opened & File access detected & Office email received & Validated \\
JSON opened & File access detected & No automatic request & Validated behavior \\
AWS ENV opened & File access detected & No activation expected & Validated behavior \\
AWS keys used & Not required on remote host & AWS token email expected, delayed & Manual test to document \\
Modify/delete & Rules 100102/100103 & Not applicable & Supported; repeat for final screenshots \\
\bottomrule
\end{tabular}
\end{table}

\section{Conclusion}

The implementation satisfies the core requirements and demonstrates a complete
local detection path. Remote activation remains format-dependent, which is a
property of the sensors rather than a reason to invent uniform behavior.
